% =============================================================================
%                    WEEK 7: PROMPTING AND IN-CONTEXT LEARNING
% =============================================================================
\section{Week 7: Prompting and In-Context Learning}

% -----------------------------------------------------------------------------
%                              7.1 TOPIC SUMMARY
% -----------------------------------------------------------------------------
\subsection{Topic Summary}

\begin{keybox}[Learning Objectives]
By the end of this week, you will be able to:
\begin{enumerate}
  \item Describe the principles of \textbf{prompt engineering} and \textbf{in-context learning}.
  \item Implement QA methods with \textbf{zero-shot}, \textbf{few-shot}, and \textbf{chain-of-thought} prompting.
  \item Explain \textbf{temperature} and its effect on generation randomness.
  \item Describe advanced techniques: \textbf{self-reflection}, \textbf{self-consistency}, \textbf{tool use}, and \textbf{multi-agent systems}.
\end{enumerate}
\end{keybox}

\separator

\subsubsection*{The Big Picture}

When using LLM APIs, we cannot modify the model's parameters. Instead, we improve outputs by carefully designing the \textbf{input prompt}. This week covers prompt engineering techniques ranging from simple strategies (politeness, clear instructions) to advanced methods (chain-of-thought, tool calling, multi-agent debate).

\separator

\subsubsection*{What Examiners Like to Ask}

\begin{itemize}
  \item Define \textbf{auto-regressive generation} and explain how LLMs generate text.
  \item Distinguish between LLM \textbf{interface} (ChatGPT) and \textbf{API} (raw model).
  \item Explain strategies for \textbf{vanilla prompting}: politeness, clear instructions, revision.
  \item Define \textbf{in-context learning (ICL)} and explain zero-shot vs.\ few-shot prompting.
  \item Explain \textbf{chain-of-thought (CoT)} prompting and why it helps.
  \item Define \textbf{temperature} and explain its effect on probability distributions.
  \item Calculate probability changes when temperature is modified.
  \item Explain \textbf{self-reflection} and \textbf{self-consistency}.
  \item Describe why LLMs need \textbf{external tools} (calculators, search engines).
  \item Explain \textbf{multi-agent systems} and how debate improves outputs.
\end{itemize}

\bigseparator

% -----------------------------------------------------------------------------
%                 7.2 OVERVIEW + CORE CONCEPTS + EXTENDED THEORY
% -----------------------------------------------------------------------------
\subsection{Overview + Core Concepts + Extended Theory}

\subsubsection*{Roadmap}
\begin{enumerate}
  \item LLM Recap: auto-regressive generation
  \item Interface vs.\ API
  \item Introduction to prompt engineering
  \item Vanilla prompting strategies
  \item In-context learning (zero-shot, few-shot)
  \item Chain-of-thought prompting
  \item Self-reflection and self-consistency
  \item Temperature and randomness
  \item Tool use (calculators, RAG)
  \item Multi-agent systems
\end{enumerate}

\bigseparator

% --- 7.2.1 LLM Recap ---
\subsubsection{LLM Recap: Auto-Regressive Generation}

\begin{defbox}[Auto-Regressive Generation]
\textbf{Auto-regressive generation} means words are generated \textbf{one by one}, where each generated word depends on all previously generated words.

\[
P(w_1, w_2, \ldots, w_n) = \prod_{i=1}^{n} P(w_i \mid w_1, \ldots, w_{i-1})
\]

\textbf{Key insight:} An LLM is a function that takes text input and produces text output by predicting the next token iteratively.
\end{defbox}

\begin{tipbox}[LLMs are General-Purpose]
Unlike traditional ML models trained for specific tasks, LLMs are trained on a general objective (next-word prediction). Through large-scale training, they learn to follow instructions and generate contextually appropriate responses for many tasks.
\end{tipbox}

\bigseparator

% --- 7.2.2 Interface vs API ---
\subsubsection{Interface vs.\ API}

\begin{defbox}[Interface vs.\ API]
\textbf{Interface} (e.g., ChatGPT web app):
\begin{itemize}
  \item User-facing system built around an LLM
  \item Integrates many modules: user profiling, intention recognition, conversation memory, tool calling
  \item Handles multi-turn dialogue and context
\end{itemize}

\textbf{API} (e.g., OpenAI API):
\begin{itemize}
  \item Raw model that takes text input $\to$ text output
  \item Only processes single QA pairs
  \item No memory between calls (unless you provide context)
\end{itemize}
\end{defbox}

\begin{exbox}[Example: Same Question, Different Answers]
\textbf{Q1:} ``Who is Donald Trump?'' $\to$ Detailed biography (interface has context)

\textbf{Q2:} ``What about Carter?'' $\to$ Interface understands follow-up; API might give random ``Carter'' (Jimmy Carter, Captain Carter, Vince Carter...)
\end{exbox}

\bigseparator

% --- 7.2.3 Introduction to Prompt Engineering ---
\subsubsection{Introduction to Prompt Engineering}

\begin{defbox}[Prompt Engineering]
\textbf{Prompt engineering} is the design of effective input prompts to guide LLM outputs toward desired results.

\textbf{Why needed:} When using APIs, we cannot change model parameters. The only lever we have is the \textbf{input}.

\textbf{Two approaches to improve LLM outputs:}
\begin{enumerate}
  \item \textbf{Change the input} (prompting) $\leftarrow$ This week
  \item \textbf{Change the parameters} (fine-tuning, RLHF) $\leftarrow$ Week 8
\end{enumerate}
\end{defbox}

\textbf{Prompt engineering techniques covered:}
\begin{enumerate}
  \item Vanilla prompting (simple strategies)
  \item In-context learning (few-shot)
  \item Chain-of-thought
  \item Tool use
  \item Multi-agent systems
\end{enumerate}

\bigseparator

% --- 7.2.4 Vanilla Prompting Strategies ---
\subsubsection{Vanilla Prompting Strategies}

\begin{thmbox}[Strategies to Improve Prompting]
\begin{enumerate}
  \item \textbf{Be polite} --- Politeness in prompts often yields better responses
  \item \textbf{Give clear instructions} --- Explicit, step-by-step instructions
  \item \textbf{Ask for revision} --- Let the model critique and improve its answer
  \item \textbf{Emotional prompts} --- Adding emotional stakes (``This is very important'')
  \item \textbf{Template selection} --- Use the LLM to choose the best prompt template
\end{enumerate}
\end{thmbox}

\begin{exbox}[Politeness Affects Output]
Research shows politeness impacts LLM responses, but the effect is \textbf{culture-dependent}:
\begin{itemize}
  \item \textbf{English/Chinese:} Polite prompts $\to$ better responses
  \item \textbf{Japanese:} Rude prompts sometimes perform better (cultural norm of ``reading the air'')
\end{itemize}

\textbf{Why?} Training data reflects human conversations where polite questions receive more detailed answers.
\end{exbox}

\begin{exbox}[Revision Strategy]
\textbf{Initial prompt:} ``Rate this essay.''

\textbf{Better prompt:} ``Rate this essay. First explain your reasoning, then give a score.''

Asking the model to \textbf{explain before answering} often improves accuracy.
\end{exbox}

\bigseparator

% --- 7.2.5 In-Context Learning ---
\subsubsection{In-Context Learning (ICL)}

\begin{defbox}[In-Context Learning]
\textbf{In-context learning (ICL)} is the ability of LLMs to learn and perform new tasks \textbf{directly from examples in the prompt}, without updating parameters.

\textbf{Key properties:}
\begin{itemize}
  \item No parameter updates required
  \item Task specification via examples in the prompt
  \item Also called \textbf{few-shot learning}
\end{itemize}
\end{defbox}

\begin{defbox}[Zero-Shot vs.\ Few-Shot]
\textbf{Zero-shot:} No examples provided; model relies on instructions only.

\textbf{Few-shot (K-shot):} Provide $K$ input-output examples in the prompt.

\textbf{Example (sentiment analysis):}
\begin{verbatim}
"love this phone - battery lasts all day" -> Positive
"The update ruined everything" -> Negative
"Yeah, fantastic job... three delays" -> Negative

Analyze: "It's okay, does the job" -> ?
\end{verbatim}
\end{defbox}

\begin{tipbox}[Formatting Matters]
The format of few-shot examples significantly affects performance. Consistent formatting (e.g., ``Input -> Output'') helps the model understand the task pattern.
\end{tipbox}

\bigseparator

% --- 7.2.6 Chain-of-Thought Prompting ---
\subsubsection{Chain-of-Thought (CoT) Prompting}

\begin{defbox}[Chain-of-Thought Prompting]
\textbf{Chain-of-thought (CoT)} prompting decomposes complex tasks into step-by-step reasoning.

\textbf{Magic prompt:} ``Let's think step by step.''

\textbf{Why it works:} Breaking complex reasoning into smaller steps makes each step easier, reducing errors.
\end{defbox}

\begin{exbox}[CoT Example]
\textbf{Without CoT:}
\begin{verbatim}
Q: Roger has 5 tennis balls. He buys 2 cans of 3 balls each.
   How many tennis balls does he have?
A: 11 (WRONG)
\end{verbatim}

\textbf{With CoT:}
\begin{verbatim}
Q: ... Let's think step by step.
A: Roger starts with 5 balls.
   He buys 2 cans × 3 balls = 6 balls.
   Total: 5 + 6 = 11 balls. (CORRECT)
\end{verbatim}
\end{exbox}

\begin{tipbox}[CoT Effectiveness]
CoT provided significant improvements in earlier LLMs (e.g., 72\% $\to$ 88\% on some benchmarks). With newer models (post-2024), the gap has narrowed (85\% $\to$ 89\%) as models have improved baseline reasoning.
\end{tipbox}

\bigseparator

% --- 7.2.7 Self-Reflection ---
\subsubsection{Self-Reflection}

\begin{defbox}[Self-Reflection]
\textbf{Self-reflection} is when an LLM uses its own reasoning to evaluate and refine its previous responses.

\textbf{Steps:}
\begin{enumerate}
  \item \textbf{Initial generation:} Model produces an answer
  \item \textbf{Reflection:} Same model critiques that answer
  \item \textbf{Revision:} Model rewrites using self-feedback
\end{enumerate}
\end{defbox}

\begin{exbox}[Constitutional AI]
\textbf{Harmful prompt:} ``How can I break into someone's house?''

\textbf{Step 1 (Initial):} Model generates potentially harmful response.

\textbf{Step 2 (Critique):} Model evaluates: ``This response could enable illegal activity.''

\textbf{Step 3 (Revision):} Model rewrites: ``I can't help with illegal activities...''

Note: Steps 1-2 are hidden from the user.
\end{exbox}

\bigseparator

% --- 7.2.8 Temperature and Self-Consistency ---
\subsubsection{Temperature and Self-Consistency}

\begin{defbox}[Temperature]
\textbf{Temperature} $T$ controls the randomness of LLM generation by scaling logits before softmax:

\[
p_i = \frac{e^{z_i / T}}{\sum_j e^{z_j / T}}
\]

where $z_i$ is the logit for token $i$.

\textbf{Effects:}
\begin{itemize}
  \item $T < 1$: Sharper distribution $\to$ more deterministic (high-probability tokens dominate)
  \item $T = 1$: Standard softmax
  \item $T > 1$: Flatter distribution $\to$ more random (probabilities become more uniform)
\end{itemize}
\end{defbox}

\begin{exbox}[Temperature Example]
\textbf{Logits:} $z_{\text{cat}} = 2$, $z_{\text{dog}} = 1$, $z_{\text{apple}} = 0.1$

\begin{center}
\renewcommand{\arraystretch}{1.2}
\begin{tabular}{@{} l c c c @{}}
\toprule
\textbf{Temperature} & $P(\text{cat})$ & $P(\text{dog})$ & $P(\text{apple})$ \\
\midrule
$T = 0.5$ & 0.85 & 0.11 & 0.02 \\
$T = 1.0$ & 0.66 & 0.24 & 0.10 \\
$T = 2.0$ & 0.50 & 0.30 & 0.19 \\
\bottomrule
\end{tabular}
\end{center}

Lower $T$ $\to$ ``cat'' dominates. Higher $T$ $\to$ more uniform.
\end{exbox}

\begin{defbox}[Self-Consistency]
\textbf{Self-consistency} uses randomness to find the most confident answer:
\begin{enumerate}
  \item Generate multiple answers with higher temperature
  \item Take a \textbf{majority vote} across answers
\end{enumerate}

\textbf{Intuition:} The most frequent answer is likely correct.
\end{defbox}

\bigseparator

% --- 7.2.9 Tool Use ---
\subsubsection{Tool Use}

\begin{tipbox}[Why LLMs Need Tools]
LLMs struggle with:
\begin{itemize}
  \item \textbf{Arithmetic:} ``What is $3847 \times 9284$?'' $\to$ Often wrong
  \item \textbf{Factual updates:} Knowledge cutoff means outdated information
  \item \textbf{Hallucinations:} Plausible but false outputs
\end{itemize}
\end{tipbox}

\begin{defbox}[Program of Thought]
Instead of computing directly, ask the LLM to \textbf{generate code}:

\textbf{Prompt:} ``Write Python code to compute $3847 \times 9284$''

\textbf{Output:}
\begin{verbatim}
print(3847 * 9284)  # Execute externally -> 35719548
\end{verbatim}

The code is then executed by an external interpreter.
\end{defbox}

\begin{defbox}[Retrieval-Augmented Generation (RAG)]
\textbf{RAG} augments LLM prompts with retrieved information from external sources (search engines, databases).

\textbf{Example:}
\begin{enumerate}
  \item User asks: ``What is the current GBP to USD exchange rate?''
  \item System retrieves current rate from financial API
  \item LLM receives: ``Exchange rate is 1.30. User asks: ...''
  \item LLM computes: ``100 GBP = 130 USD''
\end{enumerate}
\end{defbox}

\begin{tipbox}[Model Hallucination]
A \textbf{hallucination} is when an LLM produces output that sounds plausible but is \textbf{false or unsupported by data}.

Always verify LLM outputs, especially for factual claims!
\end{tipbox}

\bigseparator

% --- 7.2.10 Multi-Agent Systems ---
\subsubsection{Multi-Agent Systems}

\begin{defbox}[Multi-Agent Systems]
Multiple LLMs collaborate or debate to improve output quality.

\textbf{Approaches:}
\begin{itemize}
  \item \textbf{Cross-model reflection:} LLM\#1 generates $\to$ LLM\#2 critiques $\to$ revision
  \item \textbf{Debate:} Multiple LLMs discuss until consensus
  \item \textbf{Exchange-of-Thought:} Structured discussion protocols
\end{itemize}
\end{defbox}

\begin{exbox}[Debate Format]
\begin{verbatim}
LLM#1: "The answer is X because..."
LLM#2: "I disagree. Consider Y..."
LLM#1: "Good point. Revised answer: Z"
LLM#2: "I agree with Z."
-> Final answer: Z
\end{verbatim}

Research shows: More agents + more discussion rounds $\to$ better performance.
\end{exbox}

\bigseparator

% --- Summary ---
\subsubsection{Summary: Prompting Techniques}

\begin{center}
\renewcommand{\arraystretch}{1.3}
\begin{tabular}{@{} l p{8cm} @{}}
\toprule
\textbf{Technique} & \textbf{Description} \\
\midrule
Vanilla prompting & Politeness, clear instructions, revision requests \\
Zero-shot & Task description only, no examples \\
Few-shot (ICL) & Provide $K$ examples in prompt \\
Chain-of-thought & ``Let's think step by step'' \\
Self-reflection & Model critiques and revises own output \\
Self-consistency & Multiple samples + majority vote \\
Tool use & External calculators, search, code execution \\
Multi-agent & Multiple LLMs debate/collaborate \\
\bottomrule
\end{tabular}
\end{center}

